\documentclass[12pt]{article}
\usepackage[english]{babel}
\usepackage{natbib}
\usepackage{url}
\usepackage[utf8x]{inputenc}
\usepackage{amsmath}
\usepackage{graphicx}
\graphicspath{{images/}}
\usepackage{parskip}
\usepackage{fancyhdr}
\usepackage{vmargin}
\setmarginsrb{3 cm}{2.5 cm}{3 cm}{2.5 cm}{1 cm}{1.5 cm}{1 cm}{1.5 cm}

\usepackage{eurosym}                %simbolo dell'euro
\usepackage{listings}
\usepackage[usenames,dvipsnames,svgnames,table]{xcolor}
% Imposta lo spazio nella list of listing in modo simile alla list of figures/tables
%\makeatletter
%\let\my@chapter\@chapter
%\renewcommand*{\@chapter}{%
%  \addtocontents{lol}{\protect\addvspace{10pt}}%
%  \my@chapter}
%\makeatother


\definecolor{codegreen}{rgb}{0,0.6,0}
\definecolor{codegray}{rgb}{0.5,0.5,0.5}
\definecolor{backcolor}{rgb}{0.98,0.98,0.98}

\renewcommand{\lstlistingname}{Codice}% Listing -> codice
\renewcommand{\lstlistlistingname}{Elenco dei frammenti di codice}% List of Listings -> Frammenti di codice

\lstdefinestyle{mystyle}{
    backgroundcolor=\color{backcolor},   
    commentstyle=\color{Peach}\ttfamily,
    keywordstyle=\color{RoyalBlue},
    numberstyle=\tiny\color{codegray},
    stringstyle=\color{SeaGreen}\ttfamily,
    basicstyle=\footnotesize\ttfamily,
    breakatwhitespace=false,         
    breaklines=true,                 
    captionpos=b,                    
    keepspaces=true,                 
    numbers=left,                    
    numbersep=5pt,                  
    showspaces=false,                
    showstringspaces=false,
    showtabs=false,                  
    tabsize=2,
    frame=trbl, % draw a frame at the top, right, left and bottom of the listing
	frameround=ftff, % angolo in basso a destro curvo
	framesep=4pt, % quarter circle size of the round corners,
	inputencoding=utf8,
    extendedchars=true,
    literate={á}{{\'a}}1 {à}{{\`a}}1 {é}{{\'e}}1 {è}{{\`e}}1 {ù}{{\`u}}1 {ò}{{\`o}}1 {ì}{{\`i}}1,
    belowskip=1em,
    aboveskip=1em,
}

 
\lstset{style=mystyle}

\lstdefinelanguage{JavaScript}
{
  % list of keywords
  morekeywords={ true, false, catch, function, break,	new, class, extends, var, require, switch, return, import, if, while, for, this, View, Text, StyleSheet},
  sensitive=false, % keywords are not case-sensitive
  morecomment=[l]{//}, % l is for line comment
  morecomment=[s]{/*}{*/}, % s is for start and end delimiter
  morestring=[b]' % defines that strings are enclosed in double quotes
}

\lstdefinelanguage{JSON}
{
  % list of keywords
  morekeywords={string, boolean, int, Array, Node, Asset, AssetDetail, Filter, FilterItem},
  sensitive=false, % keywords are not case-sensitive
  morecomment=[l]{//}, % l is for line comment
  morecomment=[s]{/*}{*/}, % s is for start and end delimiter
  morestring=[b]" % defines that strings are enclosed in double quotes
}

\lstdefinelanguage{URM}
{
	% list of keywords
	morekeywords={ S, J, T, Z, I},
	sensitive=false, % keywords are not case-sensitive
	morecomment=[l]{//}, % l is for line comment
	morecomment=[s]{/*}{*/}, % s is for start and end delimiter
	morestring=[b]' % defines that strings are enclosed in double quotes
}

\lstdefinelanguage{RDFA}{
	language=html,
	sensitive=true, 
	alsoletter={<>=-},
	ndkeywords={
		% General
		=,
		% HTML attributes
		charset=, id=, width=, height=, property=, about=, rel=, rev=, prefix=, vocab=, content=, datatype=
	},  
	morecomment=[s]{<!--}{-->},
	tag=[s]
}

%\tightlist per compatibilità con pandoc
\providecommand{\tightlist}{%
	\setlength{\itemsep}{0pt}\setlength{\parskip}{0pt}}


\usepackage[labelfont=bf]{caption}

\usepackage[protrusion=true,expansion=true]{microtype} % Better typography
\usepackage{graphicx} % Required for including pictures
\usepackage{wrapfig} % Allows in-line images


\usepackage{subfig}
\usepackage{hyperref}
\usepackage{placeins}
\usepackage{sourcecodepro}
\usepackage{hyperref}                   % collegamenti ipertestuali

\usepackage[colorinlistoftodos,prependcaption]{todonotes} %todo

\usepackage{amsmath}
\usepackage{mathtools}

\usepackage{float}
\usepackage{algorithm}
\usepackage{algpseudocode} % https://en.wikibooks.org/wiki/LaTeX/Algorithms#Typesetting_using_the_algorithmicx_package
\usepackage{amsmath} 
\usepackage{amssymb}  %$\mathbb{N}$ per il simbolo dei numeri naturali 

\usepackage{enumerate} % permette di personalizzare enumerate
\usepackage{enumitem}
\usepackage{xmpincl}	%Aggiunge metadati sulla licenza CC
\usepackage{xspace}

\usepackage{bussproofs}

\usepackage{listings}
\usepackage{color}

\definecolor{dkgreen}{rgb}{0,0.6,0}
\definecolor{gray}{rgb}{0.5,0.5,0.5}
\definecolor{mauve}{rgb}{0.58,0,0.82}

\lstset{frame=tb,
  language=Java,
  aboveskip=3mm,
  belowskip=3mm,
  showstringspaces=false,
  columns=flexible,
  basicstyle={\small\ttfamily},
  numbers=none,
  numberstyle=\tiny\color{gray},
  keywordstyle=\color{blue},
  commentstyle=\color{dkgreen},
  stringstyle=\color{mauve},
  breaklines=true,
  breakatwhitespace=true,
  tabsize=3
}


\makeatletter
\renewcommand\@biblabel[1]{\textbf{#1.}} % Change the square brackets for each bibliography item from '[1]' to '1.'
\renewcommand{\@listI}{\itemsep=0pt} % Reduce the space between items in the itemize and enumerate environments and the bibliography

\renewcommand{\maketitle}{ % Customize the title - do not edit title and author name here, see the TITLE block below
	\begin{flushright} % Right align
		{\LARGE\@title} % Increase the font size of the title
		
		\vspace{50pt} % Some vertical space between the title and author name
		
		{\large\@author} % Author name
		\\\@date % Date
		
		\vspace{100pt} % Some vertical space between the author block and abstract
	\end{flushright}
}

%% breakablealgorithm http://tex.stackexchange.com/questions/33866/algorithm-tag-and-page-break
\makeatletter
\newenvironment{breakablealgorithm}
{% \begin{breakablealgorithm}
	\begin{center}
		\refstepcounter{algorithm}% New algorithm
		\hrule height.8pt depth0pt \kern2pt% \@fs@pre for \@fs@ruled
		\renewcommand{\caption}[2][\relax]{% Make a new \caption
			{\raggedright\textbf{\ALG@name~\thealgorithm} ##2\par}%
			\ifx\relax##1\relax % #1 is \relax
			\addcontentsline{loa}{algorithm}{\protect\numberline{\thealgorithm}##2}%
			\else % #1 is not \relax
			\addcontentsline{loa}{algorithm}{\protect\numberline{\thealgorithm}##1}%
			\fi
			\kern2pt\hrule\kern2pt
		}
	}{% \end{breakablealgorithm}
	\kern2pt\hrule\relax% \@fs@post for \@fs@ruled
\end{center}
}
\makeatother

\makeatletter % trattino con punto sopra
\newcommand{\dotminus}{\mathbin{\text{\@dotminus}}}

\newcommand{\@dotminus}{%
	\ooalign{\hidewidth\raise1ex\hbox{.}\hidewidth\cr$\m@th-$\cr}%
}
\makeatother

\DeclarePairedDelimiter{\ceil}{\lceil}{\rceil}
\DeclarePairedDelimiter{\floor}{\lfloor}{\rfloor}


\let\oldtext\text
\renewcommand{\text}[1]{\oldtext{\normalfont\sffamily #1}}
\newcommand{\false}{\text{ false }}
\newcommand{\true}{\text{ true }}
\newcommand{\fn}{\text{fn }}
\newcommand{\iif}{\text{ if }}
\newcommand{\then}{\text{ then }}
\newcommand{\eelse}{\text{ else }}
\newcommand{\Nat}{\text{ Nat }}
\newcommand{\Bool}{\text{ Bool }}

\newcommand{\IF}[3]{\ensuremath{\text{if}\ #1\ \text{then}\ #2\ \text{else}\ #3}}

\newcommand{\vbar}{\:|\:}
\newcommand{\proj}{.\_}

\newcommand{\throw}[1]{\ensuremath{\text{throw} \: #1}}
\newcommand{\trycatch}[2]{\ensuremath{\text{try} \: #1 \: \text{catch} \: #2}}
\newcommand{\TExc}{\ensuremath{T_{exc}}}
\newcommand{\myrule}[1]{\textsc{(#1)}}

\newcommand{\class}[3]{\text{class} \: #1 \: \text{extends} \: #2 \: \{\: #3 \:\}}
\newcommand{\construct}[3]{\ensuremath{#1(#2)\{\: #3\:\}}}
\newcommand{\this}{\text{this}}
\newcommand{\super}{\text{super}}

\newcommand{\method}[4]{\ensuremath{#1 \: #2(#3) \: \{ \: \text{return} \: #4;\: \} }}
\newcommand{\new}[2]{\ensuremath{\text{new}\: #1(  #2  )}}

\newcommand{\stype}{\ensuremath{<:}}
\newcommand{\fields}[1]{fields(#1)}
\newcommand{\mbody}[1]{mbody(#1)}
\newcommand{\mtype}[1]{mtype(#1)}
\newcommand{\IN}[1]{\text{in} \: #1}


% Shorcut per prooftree
\newcommand{\LL}[1]{\LeftLabel{\myrule{#1}}}
\newcommand{\AC}[1]{\AxiomC{#1}}
\newenvironment{bprooftree}
{\leavevmode\hbox\bgroup}
{\DisplayProof\egroup}

\title{Esercizi}                                % Title
\author{AALP}                               % Author
\date{3 Dec 2016}                                         % Date

\makeatletter
\let\thetitle\@title
\let\theauthor\@author
\let\thedate\@date
\makeatother

\pagestyle{fancy}
\fancyhf{}
\rhead{\theauthor}
\lhead{\thetitle}
\cfoot{\thepage}

\begin{document}

%%%%%%%%%%%%%%%%%%%%%%%%%%%%%%%%%%%%%%%%%%%%%%%%%%%%%%%%%%%%%%%%%%%%%%%%%%%%%%%%%%%%%%%%%

\begin{titlepage}
    \centering
    \vspace*{0.5 cm}
    \includegraphics[scale = 0.75]{logo.png}\\[1.0 cm]  % University Logo
    \textsc{\LARGE Aspetti Avanzati dei Linguaggi \newline\newline di Programmazione}\\[2.0 cm]  % University Name
    %\textsc{\Large CSX-325}\\[0.5 cm]               % Course Code
    %\rule{\linewidth}{0.2 mm} \\[0.4 cm]
    { \huge \bfseries \thetitle}\\
    \rule{\linewidth}{0.2 mm} \\[1.5 cm]
    
    \begin{minipage}{0.4\textwidth}
        \begin{flushleft} \large
            \emph{Autori:}\\
            Campese Stefano\\
            Enrico Savoca\\
        \end{flushleft}
    \end{minipage}~
    
    
    
    
    
    
    
\end{titlepage}

%%%%%%%%%%%%%%%%%%%%%%%%%%%%%%%%%%%%%%%%%%%%%%%%%%%%%%%%%%%%%%%%%%%%%%%%%%%%%%%%%%%%%%%%%

\tableofcontents
\pagebreak

%%%%%%%%%%%%%%%%%%%%%%%%%%%%%%%%%%%%%%%%%%%%%%%%%%%%%%%%%%%%%%%%%%%%%%%%%%%%%%%%%%%%%%%%%
  

\subsection{Il mini linguaggio funzionale (note 2)}
\subsubsection*{Esercizio 1.1}
La relazione di riduzione data definisce una strategia efficiente per la valutazione del ter-
mine if-then-else, che permette cioè di valutare unicamente il ramo scelto dalla valutazione della guardia
booleana. Ridefinire la semantica operazionale del linguaggio in modo che adotti una strategia non efficiente per il costrutto if-then-else, valutando entambi i rami del costrutto condizionale.

\subparagraph{Svolgimento}
Per ottenere una semantica meno efficiente, si cambiano gli assiomi (IF-TRUE) e (IF-FALSE) e si aggiungono le regole (THEN) ed (ELSE). (IF) si mantiene tale.

\subparagraph{Assiomi}
\begin{prooftree}
	\AxiomC{}
	\LeftLabel{\textsc{(IF-TRUE)}}
	\UnaryInfC{$if \: TRUE \: then \: v_1 \: else \: v_2 \rightarrow \: v_1$}
\end{prooftree}

\begin{prooftree}
	\AxiomC{}
	\LeftLabel{\textsc{(IF-FALSE)}}
	\UnaryInfC{$if \: FALSE \: then \: v_1 \: else \: v_2 \rightarrow \: v_2$}
\end{prooftree}

\subparagraph{Regole}
\begin{prooftree}
	\AxiomC{$M_2 \rightarrow \: M_2'$}
	\LeftLabel{\textsc{(THEN)}}
	\UnaryInfC{$if \: v_1 \: then \: M_2 \: else \: M_3 \rightarrow \: 
	if \: v_1 \: then \: M_2' \: else \: M_3$}
\end{prooftree}

\begin{prooftree}
	\AxiomC{$M_3 \rightarrow \: M_3'$}
	\LeftLabel{\textsc{(ELSE)}}
	\UnaryInfC{$if \: v_1 \: then \: v_2 \: else \: M_3 \rightarrow \: 
	if \: v_1 \: then \: v_2 \: else \: M_3'$} 
\end{prooftree} 

\subsubsection*{Esercizio 1.4}
La valutazione, cioè l’esecuzione del programma, è deterministica. Se $M \rightarrow$ M' e
$M \rightarrow$ M'' allora M'=M''.\\
Dimostrare la proposizione precedente per induzione sulla struttura del termine M.

\subparagraph{Dimostrazione}
Si dimostra per induzione sulla derivazione.

Casi base:
\begin{itemize}
	\item M=true \\
	\indent	Se true $\rightarrow$ M' e true $\rightarrow$ M'' $\implies$ M' = M''\\
	\indent Vero perchè true è un valore finale;
	
 \item M=false \\
	\indent	Se false $\rightarrow$ M' e false $\rightarrow$ M'' $\implies$ M' = M''\\
	\indent Vero perchè false è un valore finale;

\item M=fn x.M \\
	\indent	Se fn x.M $\rightarrow$ M' e fn x.M $\rightarrow$ M'' $\implies$ M' = M''\\
	\indent Vero perchè fn x.M è un valore finale;
	\item M=n \\
	\indent	Se n $\rightarrow$ M' e n $\rightarrow$ M'' $\implies$ M' = M''\\
	\indent Vero perchè n è un valore finale.
\end{itemize}
 
Passo induttivo:
\begin{itemize}
	\item Tesi: $M=M_1+M_2 \:\rightarrow\: M', \:M_1+M_2 \:\rightarrow$ M'' $\implies$ M'=M''\\
	\indent Per ipotesi induttiva: 
	\begin{itemize}
		\item Se $M_1 \:\rightarrow \: M_1'$ e $M_1 \:\rightarrow \: M_1''$ allora $M_1'=M_1''$;
		\item Se $M_2 \:\rightarrow \: M_2'$ e $M_2 \:\rightarrow \: M_2''$ allora $M_2'=M_2''$.
	\end{itemize} 
	3 Casi possibili:
	\begin{enumerate}[label=\alph*)]
		\item Derivazione tramite l'uso della regola (SUM-LEFT): 
		\begin{prooftree}
					\AxiomC{$M_1\:\rightarrow\:M_1'$}
					\LeftLabel{\textsc{(SUM-LEFT)}}
					\UnaryInfC{$M_1+M_2 \:\rightarrow \: M_1'+M_2=M'$}
				\end{prooftree} 
				\indent In questo caso $M_1'=M_1''$ per ipotesi 
						induttiva e quindi anche $M'=M''$;\\
				\indent La derivazione risulta deterministica applicando questa regola.
		\item Derivazione tramite l'uso della regola (SUM-RIGHT): 
		\begin{prooftree}
					\AxiomC{$M_2\:\rightarrow\:M_2'$}
					\LeftLabel{\textsc{(SUM-RIGHT)}}
					\UnaryInfC{$v_1+M_2 \:\rightarrow \: v_1+M_2'=M'$}
				\end{prooftree} 
				\indent In questo caso $M_2'=M_2''$ per ipotesi 
						induttiva e quindi anche $M'=M''$;\\
				\indent La derivazione risulta deterministica applicando questa regola.
		\item Derivazione tramite l'uso della regola (SUM): 
		\begin{prooftree}
					\AxiomC{con n'=$n_1+n_2$}
					\LeftLabel{\textsc{(SUM)}}
					\UnaryInfC{$n_1+n_2 \:\rightarrow \: n'$}
				\end{prooftree} 
				\indent Il risultato della somma è univoco e di conseguenza la derivazione è deterministica. 
	\end{enumerate}
	In tutti e 3 i casi, i soli possibili, la derivazione è deterministica, quindi la tesi 
	è dimostrata.\\
	
	
	\item Tesi: $M=M_1-M_2 \:\rightarrow\: M', \:M_1-M_2 \:\rightarrow$ M'' $\implies$ M'=M''\\
	\indent Per ipotesi induttiva: 
	\begin{itemize}
		\item Se $M_1 \:\rightarrow \: M_1'$ e $M_1 \:\rightarrow \: M_1''$ allora $M_1'=M_1''$;
		\item Se $M_2 \:\rightarrow \: M_2'$ e $M_2 \:\rightarrow \: M_2''$ allora $M_2'=M_2''$.
	\end{itemize} 
	3 Casi possibili:
	\begin{enumerate}[label=\alph*)]
		\item Derivazione tramite l'uso della regola (MINUS-LEFT): 
		\begin{prooftree}
					\AxiomC{$M_1\:\rightarrow\:M_1'$}
					\LeftLabel{\textsc{(MINUS-LEFT)}}
					\UnaryInfC{$M_1-M_2 \:\rightarrow \: M_1'-M_2=M'$}
				\end{prooftree} 
				\indent In questo caso $M_1'=M_1''$ per ipotesi 
						induttiva e quindi anche $M'=M''$;\\
				\indent La derivazione risulta deterministica applicando questa regola.
		\item Derivazione tramite l'uso della regola (MINUS-RIGHT): 
		\begin{prooftree}
					\AxiomC{$M_2\:\rightarrow\:M_2'$}
					\LeftLabel{\textsc{(MINUS-RIGHT)}}
					\UnaryInfC{$v_1-M_2 \:\rightarrow \: v_1-M_2'=M'$}
				\end{prooftree} 
				\indent In questo caso $M_2'=M_2''$ per ipotesi 
						induttiva e quindi anche $M'=M''$;\\
				\indent La derivazione risulta deterministica applicando questa regola.
		\item Derivazione tramite l'uso della regola (MINUS): 
		\begin{prooftree}
					\AxiomC{con n'=$n_1-n_2$}
					\LeftLabel{\textsc{(MINUS)}}
					\UnaryInfC{$n_1-n_2 \:\rightarrow \: n'$}
				\end{prooftree} 
				\indent Il risultato della differenza è univoco e di conseguenza la derivazione è deterministica. 
	\end{enumerate}
	In tutti e 3 i casi, i soli possibili, la derivazione è deterministica, quindi la tesi 
	è dimostrata.\\
	
	
	\item Tesi: $M=if\:M_1\:then\:M_2\:else\:M_3\:\rightarrow\: M',\:if\:M_1\:then\:M_2\:else\:M_3\:\rightarrow\: M''\:\implies$ M'=M''\\
	\indent Per ipotesi induttiva: 
	\begin{itemize}
		\item Se $M_1 \:\rightarrow \: M_1'$ e $M_1 \:\rightarrow \: M_1''$ allora $M_1'=M_1''$.
	\end{itemize}  
	3 Casi possibili:
	\begin{enumerate}[label=\alph*)]
		\item Derivazione tramite l'uso della regola (IF): 
		\begin{prooftree}
					\AxiomC{$M_1\:\rightarrow\:M_1'$}
					\LeftLabel{\textsc{(IF)}}
					\UnaryInfC{$if\:M_1\:then\:M_2\:else\:M_3\:\rightarrow\:if\:M_1'\:then\:M_2\:else\:M_3=M'$}
				\end{prooftree} 
				\indent In questo caso $M_1'=M_1''$ per ipotesi 
						induttiva e quindi anche $M'=M''$;\\
				\indent La derivazione risulta deterministica applicando questa regola.
		\item Derivazione tramite l'uso dell'assioma (IF-TRUE): 
		\begin{prooftree}
					\AxiomC{}
					\LeftLabel{\textsc{(IF-TRUE)}}
					\UnaryInfC{$if\:true\:then\:M_2\:else\:M_3\:\rightarrow\:M_2=M'$}
				\end{prooftree} 
				\indent In questo caso $M_2=M'=M''$ per applicazione dell'assioma e 
				la derivazione risulta deterministica.  
		\item Derivazione tramite l'uso della regola (IF-FALSE): 
		\begin{prooftree}
					\AxiomC{}
					\LeftLabel{\textsc{(IF-FAlSE)}}
					\UnaryInfC{$if\:false\:then\:M_2\:else\:M_3\:\rightarrow\:M_3=M'$}
				\end{prooftree} 
				\indent In questo caso $M_3=M'=M''$ per applicazione dell'assioma e 
				la derivazione risulta deterministica.  
	\end{enumerate} 
	In tutti e 3 i casi, i soli possibili, la derivazione è deterministica, quindi la tesi 
	è dimostrata.\\	
	
	
		\item Tesi: $M=M_1\:M_2 \:\rightarrow\: M', \:M_1\:M_2 \:\rightarrow$ M'' $\implies$ M'=M''\\
	\indent Per ipotesi induttiva: 
	\begin{itemize}
		\item Se $M_1 \:\rightarrow \: M_1'$ e $M_1 \:\rightarrow \: M_1''$ allora $M_1'=M_1''$;
		\item Se $M_2 \:\rightarrow \: M_2'$ e $M_2 \:\rightarrow \: M_2''$ allora $M_2'=M_2''$.
	\end{itemize}  
	2 Casi possibili:
	\begin{enumerate}[label=\alph*)]
		\item Derivazione tramite l'uso della regola (APP1): 
		\begin{prooftree}
					\AxiomC{$M_1\:\rightarrow\:M_1'$}
					\LeftLabel{\textsc{(APP1)}}
					\UnaryInfC{$M_1\:M_2 \:\rightarrow \: M_1'\:M_2=M'$}
				\end{prooftree} 
				\indent In questo caso $M_1'=M_1''$ per ipotesi 
						induttiva e quindi anche $M'=M''$;\\
				\indent La derivazione risulta deterministica applicando questa regola.
		\item Derivazione tramite l'uso della regola (APP2): 
		\begin{prooftree}
					\AxiomC{$M_2\:\rightarrow\:M_2'$}
					\LeftLabel{\textsc{(APP2)}}
					\UnaryInfC{$v_1\:M_2 \:\rightarrow \: v_1\:M_2'=M'$}
				\end{prooftree} 
				\indent In questo caso $M_2'=M_2''$ per ipotesi 
						induttiva e quindi anche $M'=M''$;\\
				\indent La derivazione risulta deterministica applicando questa regola. 
	\end{enumerate} 
	In tutti e 2 i casi, i soli possibili, la derivazione è deterministica, quindi la tesi 
	è dimostrata.\\	
	
	
		\item Tesi: $M=fn\:x.M\:v\:\rightarrow \:M{x:=v}=M'$, $fn\:x.M\:v\:\rightarrow \:M{x:=v}=M''\implies$ M'=M''\\  
 	Derivazione tramite l'uso dell'assioma (BETA): 
		\begin{prooftree}
					\AxiomC{}
					\LeftLabel{\textsc{(BETA)}}
					\UnaryInfC{$fn\:x.M\:v\:\rightarrow \:M{x:=v}=M'$}
				\end{prooftree} 
				\indent Tramite l'applicazione dell'assioma si ottiene M{x:=v}=M'=M''.\\
				\indent La derivazione risulta quindi deterministica e la tesi risulta dimostrata.\\	
\end{itemize}

La derivazione è determinstica all'applicazione di tutte le regole e gli assiomi del linguaggio, di conseguenza è deterministica l'esecuzione del programma.\\ 
\begin{flushright}
$\Box$
\end{flushright}


\subsubsection*{Esercizio 1.5}
Descrivere la valutazione del termine ((fn x.3) (fn y.y)) ((fn z.if z then 1 else 0) (false)).
Modificare le regole di valutazione in modo tale che, mantenendo una strategia call-by-value, il termine precedente evolva in un termine stuck in meno passi di riduzione. Scrivere le regole di valutazione della strategia call-by-name e valutare il termine precedente secondo questa strategia.\\

\subparagraph{Svolgimento}
Valutazione tramite strategia call-by-value: \\
\textbf{passo 1}
\begin{prooftree} 
	\AxiomC{}
	\LeftLabel{\textsc{(BETA)}}
	\UnaryInfC{$(\fn x.3)(\fn y.y)\: \rightarrow \:3$}
	\LeftLabel{\textsc{(APP1)}}
	\UnaryInfC{$((\fn x.3)(\fn y.y))\:((\fn z.if\:z\:then\:1\:else\:0)(false))\: \rightarrow 
	\:3\:((\fn z.if\:z\:then\:1\:else\:0)(false))$}
\end{prooftree}
\textbf{passo 2}
\begin{prooftree} 
	\AxiomC{}
	\LeftLabel{\textsc{(BETA)}}
	\UnaryInfC{$(\fn z.if\:z\:then\:1\:else\:0)(false)\: \rightarrow \:if\:false\:then\:1\:else\:0$}
	\LeftLabel{\textsc{(APP2)}}
	\UnaryInfC{$3\:((\fn z.if\:z\:then\:1\:else\:0)(false))\: \rightarrow 
	\:3\:(if\:false\:then\:1\:else\:0)$}
\end{prooftree}
\textbf{passo 3}
\begin{prooftree} 
	\AxiomC{}
	\LeftLabel{\textsc{(IF-FALSE)}}
	\UnaryInfC{$if\:false\:then\:1\:else\:0\: \rightarrow \:0$}
	\LeftLabel{\textsc{(APP2)}}
	\UnaryInfC{$3\:(if\:false\:then\:1\:else\:0)\: \rightarrow 
	\:3\:0$}
\end{prooftree}

3 0 costituisce uno stuck: non esistono regole o assiomi per continuare l'esecuzione di tale programma. Esecuzione terminata dopo 3 passi di derivazione. \\
Si può mantenere la strategia CBV e terminare l'esecuzione in un numero minore di passi modificando (APP2):
 
\begin{prooftree}
	\AxiomC{$N\:\rightarrow\:N'$}
	\LeftLabel{\textsc{(APP2)}}
	\UnaryInfC{$(\fn x.M)\: N\:\rightarrow\:\fn x.M\:N'$} 
\end{prooftree}

In questo modo l'esecuzione del programma terminerebbe dopo un solo passo. Infatti, ci si ritroverebbe in una forma in cui non è possibile applicare alcuna regola o assioma (3 non è una funzione).\\

Con una strategia CBN (Call-by-name) non ho (APP2) e la funzione (BETA) è la seguente:

\begin{prooftree}
	\AxiomC{$\checkmark$}
	\LeftLabel{\textsc{(BETA)}} 
	\UnaryInfC{$(\fn x.M)\: N\: \rightarrow \:M{x:=N}$}
\end{prooftree}

Con questa regola mostriamo l'esecuzione del programma:\\

\textbf{passo 1}
\begin{prooftree} 
	\AxiomC{}
	\LeftLabel{\textsc{(BETA)}}
	\UnaryInfC{$(\fn x.3)(\fn y.y)\: \rightarrow \:3$}
	\LeftLabel{\textsc{(APP1)}}
	\UnaryInfC{$((\fn x.3)(\fn y.y))\:((\fn z.if\:z\:then\:1\:else\:0)(false))\: \rightarrow 
	\:3\:((\fn z.if\:z\:then\:1\:else\:0)(false))$}
\end{prooftree}

Stuck! Non esistono regole o assiomi applicabili: APP1 non è applicabile in quanto 3 non è funzione.



\subsubsection*{Esercizio 1.6}
Consideriamo le seguenti definizioni in Scala:
def square(x:Int):Int = x*x
def sumOfSquare(x:Int,y:Int):Int = square(x)+square(y)
 Descrivere i passi di riduzione dell’espressione sumOfSquare(3,4) secondo una strategia call-
by-value analoga a quella definita nella sezione precedente. Descrivere inoltre la riduzione della
stesa espressione secondo la strategia call-by-name.
 Descrivere i passi di riduzione dell’espressione sumOfSquare(3,2+2) secondo le strategie call-by-value e call-by-name.

\subparagraph{Svolgimento}
Caso sumOfSquare(3,4).\\
CBV si comporta come CBN.\\

\textbf{passo 1}
\begin{prooftree} 
	\AxiomC{}
	\LeftLabel{\textsc{(BETA)}}
	\UnaryInfC{sumOfSquare(3, 4) $\rightarrow$ square(3) + square(4)} 
\end{prooftree}

\textbf{passo 2}
\begin{prooftree} 
	\AxiomC{}
	\LeftLabel{\textsc{(BETA)}}
	\UnaryInfC{square(3) $\rightarrow$ 3*3} 
	\LeftLabel{\textsc{(SUM-L)}}
	\UnaryInfC{square(3) + square(4) $\rightarrow$ 3*3 + square(4)} 
\end{prooftree}

\textbf{passo 3}
\begin{prooftree} 
	\AxiomC{}
	\LeftLabel{\textsc{(MULT)}}
	\UnaryInfC{3*3 $\rightarrow$ 9} 
	\LeftLabel{\textsc{(SUM-L)}}
	\UnaryInfC{3*3 + square(4) $\rightarrow$ 9 + square(4)} 
\end{prooftree}

\textbf{passo 4}
\begin{prooftree} 
	\AxiomC{}
	\LeftLabel{\textsc{(BETA)}}
	\UnaryInfC{square(4) $\rightarrow$ 4*4} 
	\LeftLabel{\textsc{(SUM-R)}}
	\UnaryInfC{9 + square(4) $\rightarrow$ 9 + 4*4} 
\end{prooftree}

\textbf{passo 5}
\begin{prooftree} 
	\AxiomC{}
	\LeftLabel{\textsc{(MULT)}}
	\UnaryInfC{4*4 $\rightarrow$ 16} 
	\LeftLabel{\textsc{(SUM-R)}}
	\UnaryInfC{9 + 4*4 $\rightarrow$ 9 + 16} 
\end{prooftree}

\textbf{passo 6}
\begin{prooftree} 
	\AxiomC{} 
	\LeftLabel{\textsc{(SUM)}}
	\UnaryInfC{9 + 16 $\rightarrow$ 25} 
\end{prooftree}

Caso sumOfSquare(3,2+2)\\
CBV:\\
\textbf{passo 1}
\begin{prooftree} 
	\AxiomC{}
	\LeftLabel{\textsc{(SUM)}}
	\UnaryInfC{2+2 $\rightarrow$ 4} 
	\LeftLabel{\textsc{(APP2)}}
	\UnaryInfC{sumOfSquare(3, 2+2) $\rightarrow$ sumOfSquare(3, 4)} 
\end{prooftree}

I passi successivi sono gli stessi del caso precedente, quindi i passi di esecuzione sono 7.\\

CBN:\\
\textbf{passo 1}
\begin{prooftree} 
	\AxiomC{} 
	\LeftLabel{\textsc{(BETA)}}
	\UnaryInfC{sumOfSquare(3, 2+2) $\rightarrow$ square(3) + square(2+2)} 
\end{prooftree}

\textbf{passo 2}
\begin{prooftree} 
	\AxiomC{} 
	\LeftLabel{\textsc{(BETA)}}
	\UnaryInfC{square(3) $\rightarrow$ 3*3}
	\LeftLabel{\textsc{(SUM-L)}}
	\UnaryInfC{square(3) + square(2+2) $\rightarrow$ 3*3 + square(2+2)} 
\end{prooftree}

\textbf{passo 3}
\begin{prooftree} 
	\AxiomC{} 
	\LeftLabel{\textsc{(MULT)}}
	\UnaryInfC{3*3 $\rightarrow$ 9}
	\LeftLabel{\textsc{(SUM-L)}}
	\UnaryInfC{3*3 + square(2+2) $\rightarrow$ 9 + square(2+2)} 
\end{prooftree}

\textbf{passo 4}
\begin{prooftree} 
	\AxiomC{} 
	\LeftLabel{\textsc{(BETA)}}
	\UnaryInfC{square(2+2) $\rightarrow$ (2+2)*(2+2)}
	\LeftLabel{\textsc{(SUM-R)}}
	\UnaryInfC{9 + square(2+2) $\rightarrow$ 9 + (2+2)*(2+2)} 
\end{prooftree}

\textbf{passo 5}
\begin{prooftree} 
	\AxiomC{} 
	\LeftLabel{\textsc{(SUM)}}
	\UnaryInfC{2+2 $\rightarrow$ 4}
	\LeftLabel{\textsc{(MULT-L)}}
	\UnaryInfC{(2+2)*(2+2) $\rightarrow$ (4)*(2+2)}
	\LeftLabel{\textsc{(SUM-R)}}
	\UnaryInfC{9 + (2+2)*(2+2) $\rightarrow$ 9 + (4)*(2+2)} 
\end{prooftree}

\textbf{passo 6}
\begin{prooftree} 
	\AxiomC{} 
	\LeftLabel{\textsc{(SUM)}}
	\UnaryInfC{2+2 $\rightarrow$ 4}
	\LeftLabel{\textsc{(MULT-R)}}
	\UnaryInfC{4*(2+2) $\rightarrow$ 4*4}
	\LeftLabel{\textsc{(SUM-R)}}
	\UnaryInfC{9 + 4*(2+2) $\rightarrow$ 9 + 4*4} 
\end{prooftree}

\textbf{passo 7}
\begin{prooftree} 
	\AxiomC{}  
	\LeftLabel{\textsc{(MULT)}}
	\UnaryInfC{4*4 $\rightarrow$ 16}
	\LeftLabel{\textsc{(SUM-R)}}
	\UnaryInfC{9 + 4*4 $\rightarrow$ 9 + 16} 
\end{prooftree}

\textbf{passo 8}
\begin{prooftree} 
	\AxiomC{}   
	\LeftLabel{\textsc{(SUM)}}
	\UnaryInfC{9 + 16 $\rightarrow$ 25} 
\end{prooftree}

I passi di esecuzione sono 8 in questo caso (CBN).

\subsubsection*{Esercizio 1.7}
Si consideri la seguente definizione in Scala:
def test(x:Int, y:Int):Int = x*x
confrontare la velocità (cioè il numero di passi) di riduzione delle seguenti espressioni secondo le strategie
CBV e CBN, indicando quale delle due è piú veloce 
\begin{enumerate}
	\item test(2,3)
	\item test(3+4,8)
	\item test(7,2*4)
	\item test(3+4,2*4)
\end{enumerate}

\subparagraph{Svolgimento}
\begin{enumerate}
	\item CBV=CBN: entrambi terminano l'esecuzione in 2 passi; 
	\item CBV<CBN: CBV termina prima perchè valuta l'espressione 3+4 subito (una sola volta). CBN invece invocherà prima la funzione, sostituendo ad x l'espressione. Essendo x presente 2 volte nella funzione, la valutazione dovrà essere effettuata 2 volte. 
	\item CBV>CBN: CBN termina prima perchè non valuta l'espressione 2*4 che non è utilizzata nella funzione test.
	\item CBV=CBN: entrambi terminano l'esecuzione in 4 passi; la differenza sta nel fatto che con CBV si svolgeranno prima le due operazioni (somma e moltiplicazione) e poi si chiamerà la funzione. Con CBN invece, si eviterà di calcolare 2*4, tuttavia si calcolerà 3+4 due volte.
\end{enumerate}



\subsubsection*{Esercizio 1.8}
Definire in Scala una funzione and che si comporta come il costrutto logico $x\&\&y$. 

\subparagraph{Svolgimento}

\begin{lstlisting}
// and logico. b1 e b2 passati by-name 
def and(b1:=> Boolean, b2:=> Boolean): Boolean =
	if(!b1) false
	else if(b2) true
		 else false; 
\end{lstlisting}













\section{Il mini linguaggio funzionale (note 2)}
\subsection*{Esercizio 1.1}
La relazione di riduzione data definisce una strategia efficiente per la valutazione del termine if-then-else, che permette cioè di valutare unicamente il ramo scelto dalla valutazione della guardia
booleana. Ridefinire la semantica operazionale del linguaggio in modo che adotti una strategia non efficiente per il costrutto if-then-else, valutando entambi i rami del costrutto condizionale.

\subsubsection*{Svolgimento}
Per ottenere una semantica meno efficiente, si cambiano gli assiomi (IF-TRUE) e (IF-FALSE) e si aggiungono le regole (THEN) ed (ELSE). (IF) si mantiene tale.

\paragraph{Assiomi}
\begin{prooftree}
	\AxiomC{}
	\LeftLabel{\textsc{(IF-TRUE)}}
	\UnaryInfC{$if \: TRUE \: then \: v_1 \: else \: v_2 \rightarrow \: v_1$}
\end{prooftree}

\begin{prooftree}
	\AxiomC{}
	\LeftLabel{\textsc{(IF-FALSE)}}
	\UnaryInfC{$if \: FALSE \: then \: v_1 \: else \: v_2 \rightarrow \: v_2$}
\end{prooftree}

\paragraph{Regole}
\begin{prooftree}
	\AxiomC{$M_2 \rightarrow \: M_2'$}
	\LeftLabel{\textsc{(THEN)}}
	\UnaryInfC{$if \: v_1 \: then \: M_2 \: else \: M_3 \rightarrow \: 
	if \: v_1 \: then \: M_2' \: else \: M_3$}
\end{prooftree}

\begin{prooftree}
	\AxiomC{$M_3 \rightarrow \: M_3'$}
	\LeftLabel{\textsc{(ELSE)}}
	\UnaryInfC{$if \: v_1 \: then \: v_2 \: else \: M_3 \rightarrow \: 
	if \: v_1 \: then \: v_2 \: else \: M_3'$} 
\end{prooftree} 

\subsection*{Esercizio 1.4}
La valutazione, cioè l’esecuzione del programma, è deterministica. Se $M \rightarrow$ M' e
$M \rightarrow$ M'' allora M'=M''.\\
Dimostrare la proposizione precedente per induzione sulla struttura del termine M.

\subsubsection*{Dimostrazione}
Si dimostra per induzione sulla derivazione.

Casi base:
\begin{itemize}
	\item M=true \\
	\indent	Se true $\rightarrow$ M' e true $\rightarrow$ M'' $\implies$ M' = M''\\
	\indent Vero perchè true è un valore finale;
	
 \item M=false \\
	\indent	Se false $\rightarrow$ M' e false $\rightarrow$ M'' $\implies$ M' = M''\\
	\indent Vero perchè false è un valore finale;

\item M=fn x.M \\
	\indent	Se fn x.M $\rightarrow$ M' e fn x.M $\rightarrow$ M'' $\implies$ M' = M''\\
	\indent Vero perchè fn x.M è un valore finale;
	\item M=n \\
	\indent	Se n $\rightarrow$ M' e n $\rightarrow$ M'' $\implies$ M' = M''\\
	\indent Vero perchè n è un valore finale.
\end{itemize}
 
Passo induttivo:
\begin{itemize}
	\item Tesi: $M=M_1+M_2 \:\rightarrow\: M', \:M_1+M_2 \:\rightarrow$ M'' $\implies$ M'=M''\\
	\indent Per ipotesi induttiva: 
	\begin{itemize}
		\item Se $M_1 \:\rightarrow \: M_1'$ e $M_1 \:\rightarrow \: M_1''$ allora $M_1'=M_1''$;
		\item Se $M_2 \:\rightarrow \: M_2'$ e $M_2 \:\rightarrow \: M_2''$ allora $M_2'=M_2''$.
	\end{itemize} 
	3 Casi possibili:
	\begin{enumerate}[label=\alph*)]
		\item Derivazione tramite l'uso della regola (SUM-LEFT): 
		\begin{prooftree}
					\AxiomC{$M_1\:\rightarrow\:M_1'$}
					\LeftLabel{\textsc{(SUM-LEFT)}}
					\UnaryInfC{$M_1+M_2 \:\rightarrow \: M_1'+M_2=M'$}
				\end{prooftree} 
				\indent In questo caso $M_1'=M_1''$ per ipotesi 
						induttiva e quindi anche $M'=M''$;\\
				\indent La derivazione risulta deterministica applicando questa regola.
		\item Derivazione tramite l'uso della regola (SUM-RIGHT): 
		\begin{prooftree}
					\AxiomC{$M_2\:\rightarrow\:M_2'$}
					\LeftLabel{\textsc{(SUM-RIGHT)}}
					\UnaryInfC{$v_1+M_2 \:\rightarrow \: v_1+M_2'=M'$}
				\end{prooftree} 
				\indent In questo caso $M_2'=M_2''$ per ipotesi 
						induttiva e quindi anche $M'=M''$;\\
				\indent La derivazione risulta deterministica applicando questa regola.
		\item Derivazione tramite l'uso della regola (SUM): 
		\begin{prooftree}
					\AxiomC{con n'=$n_1+n_2$}
					\LeftLabel{\textsc{(SUM)}}
					\UnaryInfC{$n_1+n_2 \:\rightarrow \: n'$}
				\end{prooftree} 
				\indent Il risultato della somma è univoco e di conseguenza la derivazione è deterministica. 
	\end{enumerate}
	In tutti e 3 i casi, i soli possibili, la derivazione è deterministica, quindi la tesi 
	è dimostrata.\\
	
	
	\item Tesi: $M=M_1-M_2 \:\rightarrow\: M', \:M_1-M_2 \:\rightarrow$ M'' $\implies$ M'=M''\\
	\indent Per ipotesi induttiva: 
	\begin{itemize}
		\item Se $M_1 \:\rightarrow \: M_1'$ e $M_1 \:\rightarrow \: M_1''$ allora $M_1'=M_1''$;
		\item Se $M_2 \:\rightarrow \: M_2'$ e $M_2 \:\rightarrow \: M_2''$ allora $M_2'=M_2''$.
	\end{itemize} 
	3 Casi possibili:
	\begin{enumerate}[label=\alph*)]
		\item Derivazione tramite l'uso della regola (MINUS-LEFT): 
		\begin{prooftree}
					\AxiomC{$M_1\:\rightarrow\:M_1'$}
					\LeftLabel{\textsc{(MINUS-LEFT)}}
					\UnaryInfC{$M_1-M_2 \:\rightarrow \: M_1'-M_2=M'$}
				\end{prooftree} 
				\indent In questo caso $M_1'=M_1''$ per ipotesi 
						induttiva e quindi anche $M'=M''$;\\
				\indent La derivazione risulta deterministica applicando questa regola.
		\item Derivazione tramite l'uso della regola (MINUS-RIGHT): 
		\begin{prooftree}
					\AxiomC{$M_2\:\rightarrow\:M_2'$}
					\LeftLabel{\textsc{(MINUS-RIGHT)}}
					\UnaryInfC{$v_1-M_2 \:\rightarrow \: v_1-M_2'=M'$}
				\end{prooftree} 
				\indent In questo caso $M_2'=M_2''$ per ipotesi 
						induttiva e quindi anche $M'=M''$;\\
				\indent La derivazione risulta deterministica applicando questa regola.
		\item Derivazione tramite l'uso della regola (MINUS): 
		\begin{prooftree}
					\AxiomC{con n'=$n_1-n_2$}
					\LeftLabel{\textsc{(MINUS)}}
					\UnaryInfC{$n_1-n_2 \:\rightarrow \: n'$}
				\end{prooftree} 
				\indent Il risultato della differenza è univoco e di conseguenza la derivazione è deterministica. 
	\end{enumerate}
	In tutti e 3 i casi, i soli possibili, la derivazione è deterministica, quindi la tesi 
	è dimostrata.\\
	
	
	\item Tesi: $M=if\:M_1\:then\:M_2\:else\:M_3\:\rightarrow\: M',\:if\:M_1\:then\:M_2\:else\:M_3\:\rightarrow\: M''\:\implies$ M'=M''\\
	\indent Per ipotesi induttiva: 
	\begin{itemize}
		\item Se $M_1 \:\rightarrow \: M_1'$ e $M_1 \:\rightarrow \: M_1''$ allora $M_1'=M_1''$.
	\end{itemize}  
	3 Casi possibili:
	\begin{enumerate}[label=\alph*)]
		\item Derivazione tramite l'uso della regola (IF): 
		\begin{prooftree}
					\AxiomC{$M_1\:\rightarrow\:M_1'$}
					\LeftLabel{\textsc{(IF)}}
					\UnaryInfC{$if\:M_1\:then\:M_2\:else\:M_3\:\rightarrow\:if\:M_1'\:then\:M_2\:else\:M_3=M'$}
				\end{prooftree} 
				\indent In questo caso $M_1'=M_1''$ per ipotesi 
						induttiva e quindi anche $M'=M''$;\\
				\indent La derivazione risulta deterministica applicando questa regola.
		\item Derivazione tramite l'uso dell'assioma (IF-TRUE): 
		\begin{prooftree}
					\AxiomC{}
					\LeftLabel{\textsc{(IF-TRUE)}}
					\UnaryInfC{$if\:true\:then\:M_2\:else\:M_3\:\rightarrow\:M_2=M'$}
				\end{prooftree} 
				\indent In questo caso $M_2=M'=M''$ per applicazione dell'assioma e 
				la derivazione risulta deterministica.  
		\item Derivazione tramite l'uso della regola (IF-FALSE): 
		\begin{prooftree}
					\AxiomC{}
					\LeftLabel{\textsc{(IF-FAlSE)}}
					\UnaryInfC{$if\:false\:then\:M_2\:else\:M_3\:\rightarrow\:M_3=M'$}
				\end{prooftree} 
				\indent In questo caso $M_3=M'=M''$ per applicazione dell'assioma e 
				la derivazione risulta deterministica.  
	\end{enumerate} 
	In tutti e 3 i casi, i soli possibili, la derivazione è deterministica, quindi la tesi 
	è dimostrata.\\	
	
	
		\item Tesi: $M=M_1\:M_2 \:\rightarrow\: M', \:M_1\:M_2 \:\rightarrow$ M'' $\implies$ M'=M''\\
	\indent Per ipotesi induttiva: 
	\begin{itemize}
		\item Se $M_1 \:\rightarrow \: M_1'$ e $M_1 \:\rightarrow \: M_1''$ allora $M_1'=M_1''$;
		\item Se $M_2 \:\rightarrow \: M_2'$ e $M_2 \:\rightarrow \: M_2''$ allora $M_2'=M_2''$.
	\end{itemize}  
	2 Casi possibili:
	\begin{enumerate}[label=\alph*)]
		\item Derivazione tramite l'uso della regola (APP1): 
		\begin{prooftree}
					\AxiomC{$M_1\:\rightarrow\:M_1'$}
					\LeftLabel{\textsc{(APP1)}}
					\UnaryInfC{$M_1\:M_2 \:\rightarrow \: M_1'\:M_2=M'$}
				\end{prooftree} 
				\indent In questo caso $M_1'=M_1''$ per ipotesi 
						induttiva e quindi anche $M'=M''$;\\
				\indent La derivazione risulta deterministica applicando questa regola.
		\item Derivazione tramite l'uso della regola (APP2): 
		\begin{prooftree}
					\AxiomC{$M_2\:\rightarrow\:M_2'$}
					\LeftLabel{\textsc{(APP2)}}
					\UnaryInfC{$v_1\:M_2 \:\rightarrow \: v_1\:M_2'=M'$}
				\end{prooftree} 
				\indent In questo caso $M_2'=M_2''$ per ipotesi 
						induttiva e quindi anche $M'=M''$;\\
				\indent La derivazione risulta deterministica applicando questa regola. 
	\end{enumerate} 
	In tutti e 2 i casi, i soli possibili, la derivazione è deterministica, quindi la tesi 
	è dimostrata.\\	
	
	
		\item Tesi: $M=fn\:x.M\:v\:\rightarrow \:M\{x:=v\}=M'$, $fn\:x.M\:v\:\rightarrow \:M\{x:=v\}=M''\implies$ M'=M''\\  
 	Derivazione tramite l'uso dell'assioma (BETA): 
		\begin{prooftree}
					\AxiomC{}
					\LeftLabel{\textsc{(BETA)}}
					\UnaryInfC{$fn\:x.M\:v\:\rightarrow \:M\{x:=v\}=M'$}
				\end{prooftree} 
				\indent Tramite l'applicazione dell'assioma si ottiene M{x:=v}=M'=M''.\\
				\indent La derivazione risulta quindi deterministica e la tesi risulta dimostrata.\\	
\end{itemize}

La derivazione è determinstica all'applicazione di tutte le regole e gli assiomi del linguaggio, di conseguenza è deterministica l'esecuzione del programma.\\ 
\begin{flushright}
$\Box$
\end{flushright}

\subsection*{Esercizio 1.5}
Descrivere la valutazione del termine ((fn x.3) (fn y.y)) ((fn z.if z then 1 else 0) (false)).
Modificare le regole di valutazione in modo tale che, mantenendo una strategia call-by-value, il termine precedente evolva in un termine stuck in meno passi di riduzione. Scrivere le regole di valutazione della strategia call-by-name e valutare il termine precedente secondo questa strategia.\\

\subsubsection*{Svolgimento}
Valutazione tramite strategia call-by-value: \\
\textbf{passo 1}
\begin{prooftree} 
	\AxiomC{}
	\LeftLabel{\textsc{(BETA)}}
	\UnaryInfC{$(\fn x.3)(\fn y.y)\: \rightarrow \:3$}
	\LeftLabel{\textsc{(APP1)}}
	\UnaryInfC{$((\fn x.3)(\fn y.y))\:((\fn z.if\:z\:then\:1\:else\:0)(false))\: \rightarrow 
	\:3\:((\fn z.if\:z\:then\:1\:else\:0)(false))$}
\end{prooftree}
\textbf{passo 2}
\begin{prooftree} 
	\AxiomC{}
	\LeftLabel{\textsc{(BETA)}}
	\UnaryInfC{$(\fn z.if\:z\:then\:1\:else\:0)(false)\: \rightarrow \:if\:false\:then\:1\:else\:0$}
	\LeftLabel{\textsc{(APP2)}}
	\UnaryInfC{$3\:((\fn z.if\:z\:then\:1\:else\:0)(false))\: \rightarrow 
	\:3\:(if\:false\:then\:1\:else\:0)$}
\end{prooftree}
\textbf{passo 3}
\begin{prooftree} 
	\AxiomC{}
	\LeftLabel{\textsc{(IF-FALSE)}}
	\UnaryInfC{$if\:false\:then\:1\:else\:0\: \rightarrow \:0$}
	\LeftLabel{\textsc{(APP2)}}
	\UnaryInfC{$3\:(if\:false\:then\:1\:else\:0)\: \rightarrow 
	\:3\:0$}
\end{prooftree}

3 0 costituisce uno stuck: non esistono regole o assiomi per continuare l'esecuzione di tale programma. Esecuzione terminata dopo 3 passi di derivazione. \\
Si può mantenere la strategia CBV e terminare l'esecuzione in un numero minore di passi modificando (APP2):
 
\begin{prooftree}
	\AxiomC{$N\:\rightarrow\:N'$}
	\LeftLabel{\textsc{(APP2)}}
	\UnaryInfC{$(\fn x.M)\: N\:\rightarrow\:\fn x.M\:N'$} 
\end{prooftree}

In questo modo l'esecuzione del programma terminerebbe dopo un solo passo. Infatti, ci si ritroverebbe in una forma in cui non è possibile applicare alcuna regola o assioma (3 non è una funzione).\\

Con una strategia CBN (Call-by-name) non ho (APP2) e la funzione (BETA) è la seguente:

\begin{prooftree}
	\AxiomC{$\checkmark$}
	\LeftLabel{\textsc{(BETA)}} 
	\UnaryInfC{$(\fn x.M)\: N\: \rightarrow \:M\{x:=N\}$}
\end{prooftree}

Con questa regola mostriamo l'esecuzione del programma:\\

\textbf{passo 1}
\begin{prooftree} 
	\AxiomC{}
	\LeftLabel{\textsc{(BETA)}}
	\UnaryInfC{$(\fn x.3)(\fn y.y)\: \rightarrow \:3$}
	\LeftLabel{\textsc{(APP1)}}
	\UnaryInfC{$((\fn x.3)(\fn y.y))\:((\fn z.if\:z\:then\:1\:else\:0)(false))\: \rightarrow 
	\:3\:((\fn z.if\:z\:then\:1\:else\:0)(false))$}
\end{prooftree}

Stuck! Non esistono regole o assiomi applicabili: APP1 non è applicabile in quanto 3 non è funzione.



\subsection*{Esercizio 1.6}
Consideriamo le seguenti definizioni in Scala:\\
\begin{lstlisting}
def square(x:Int):Int = x*x  
def sumOfSquare(x:Int,y:Int):Int = square(x)+square(y)
\end{lstlisting}
Descrivere i passi di riduzione dell'espressione \textbf{sumOfSquare(3,4)} secondo una strategia call-by-value analoga a quella definita nella sezione precedente.\\ 
Descrivere inoltre la riduzione della stesa espressione secondo la strategia call-by-name. \\
Descrivere i passi di riduzione dell'espressione \textbf{sumOfSquare(3,2+2)} secondo le strategie call-by-value e call-by-name.

\subsubsection*{Svolgimento}
Caso \textbf{sumOfSquare(3,4)}.\\
CBV si comporta come CBN.\\

\textbf{passo 1}
\begin{prooftree} 
	\AxiomC{}
	\LeftLabel{\textsc{(BETA)}}
	\UnaryInfC{sumOfSquare(3, 4) $\rightarrow$ square(3) + square(4)} 
\end{prooftree}

\textbf{passo 2}
\begin{prooftree} 
	\AxiomC{}
	\LeftLabel{\textsc{(BETA)}}
	\UnaryInfC{square(3) $\rightarrow$ 3*3} 
	\LeftLabel{\textsc{(SUM-L)}}
	\UnaryInfC{square(3) + square(4) $\rightarrow$ 3*3 + square(4)} 
\end{prooftree}

\textbf{passo 3}
\begin{prooftree} 
	\AxiomC{}
	\LeftLabel{\textsc{(MULT)}}
	\UnaryInfC{3*3 $\rightarrow$ 9} 
	\LeftLabel{\textsc{(SUM-L)}}
	\UnaryInfC{3*3 + square(4) $\rightarrow$ 9 + square(4)} 
\end{prooftree}

\textbf{passo 4}
\begin{prooftree} 
	\AxiomC{}
	\LeftLabel{\textsc{(BETA)}}
	\UnaryInfC{square(4) $\rightarrow$ 4*4} 
	\LeftLabel{\textsc{(SUM-R)}}
	\UnaryInfC{9 + square(4) $\rightarrow$ 9 + 4*4} 
\end{prooftree}

\textbf{passo 5}
\begin{prooftree} 
	\AxiomC{}
	\LeftLabel{\textsc{(MULT)}}
	\UnaryInfC{4*4 $\rightarrow$ 16} 
	\LeftLabel{\textsc{(SUM-R)}}
	\UnaryInfC{9 + 4*4 $\rightarrow$ 9 + 16} 
\end{prooftree}

\textbf{passo 6}
\begin{prooftree} 
	\AxiomC{} 
	\LeftLabel{\textsc{(SUM)}}
	\UnaryInfC{9 + 16 $\rightarrow$ 25} 
\end{prooftree}

Caso \textbf{sumOfSquare(3,2+2)}\\
CBV:\\
\textbf{passo 1}
\begin{prooftree} 
	\AxiomC{}
	\LeftLabel{\textsc{(SUM)}}
	\UnaryInfC{2+2 $\rightarrow$ 4} 
	\LeftLabel{\textsc{(APP2)}}
	\UnaryInfC{sumOfSquare(3, 2+2) $\rightarrow$ sumOfSquare(3, 4)} 
\end{prooftree}

I passi successivi sono gli stessi del caso precedente, quindi i passi di esecuzione sono 7.\\

CBN:\\
\textbf{passo 1}
\begin{prooftree} 
	\AxiomC{} 
	\LeftLabel{\textsc{(BETA)}}
	\UnaryInfC{sumOfSquare(3, 2+2) $\rightarrow$ square(3) + square(2+2)} 
\end{prooftree}

\textbf{passo 2}
\begin{prooftree} 
	\AxiomC{} 
	\LeftLabel{\textsc{(BETA)}}
	\UnaryInfC{square(3) $\rightarrow$ 3*3}
	\LeftLabel{\textsc{(SUM-L)}}
	\UnaryInfC{square(3) + square(2+2) $\rightarrow$ 3*3 + square(2+2)} 
\end{prooftree}

\textbf{passo 3}
\begin{prooftree} 
	\AxiomC{} 
	\LeftLabel{\textsc{(MULT)}}
	\UnaryInfC{3*3 $\rightarrow$ 9}
	\LeftLabel{\textsc{(SUM-L)}}
	\UnaryInfC{3*3 + square(2+2) $\rightarrow$ 9 + square(2+2)} 
\end{prooftree}

\textbf{passo 4}
\begin{prooftree} 
	\AxiomC{} 
	\LeftLabel{\textsc{(BETA)}}
	\UnaryInfC{square(2+2) $\rightarrow$ (2+2)*(2+2)}
	\LeftLabel{\textsc{(SUM-R)}}
	\UnaryInfC{9 + square(2+2) $\rightarrow$ 9 + (2+2)*(2+2)} 
\end{prooftree}

\textbf{passo 5}
\begin{prooftree} 
	\AxiomC{} 
	\LeftLabel{\textsc{(SUM)}}
	\UnaryInfC{2+2 $\rightarrow$ 4}
	\LeftLabel{\textsc{(MULT-L)}}
	\UnaryInfC{(2+2)*(2+2) $\rightarrow$ (4)*(2+2)}
	\LeftLabel{\textsc{(SUM-R)}}
	\UnaryInfC{9 + (2+2)*(2+2) $\rightarrow$ 9 + (4)*(2+2)} 
\end{prooftree}

\textbf{passo 6}
\begin{prooftree} 
	\AxiomC{} 
	\LeftLabel{\textsc{(SUM)}}
	\UnaryInfC{2+2 $\rightarrow$ 4}
	\LeftLabel{\textsc{(MULT-R)}}
	\UnaryInfC{4*(2+2) $\rightarrow$ 4*4}
	\LeftLabel{\textsc{(SUM-R)}}
	\UnaryInfC{9 + 4*(2+2) $\rightarrow$ 9 + 4*4} 
\end{prooftree}

\textbf{passo 7}
\begin{prooftree} 
	\AxiomC{}  
	\LeftLabel{\textsc{(MULT)}}
	\UnaryInfC{4*4 $\rightarrow$ 16}
	\LeftLabel{\textsc{(SUM-R)}}
	\UnaryInfC{9 + 4*4 $\rightarrow$ 9 + 16} 
\end{prooftree}

\textbf{passo 8}
\begin{prooftree} 
	\AxiomC{}   
	\LeftLabel{\textsc{(SUM)}}
	\UnaryInfC{9 + 16 $\rightarrow$ 25} 
\end{prooftree}

I passi di esecuzione sono 8 in questo caso (CBN).

\subsection*{Esercizio 1.7}
Si consideri la seguente definizione in Scala:

\begin{lstlisting}
def test(x:Int, y:Int):Int = x*x
\end{lstlisting}
confrontare la velocit\`a (cio\`e il numero di passi) di riduzione delle seguenti espressioni secondo le strategie CBV e CBN, indicando quale delle due \`e pi\`u veloce. 
\begin{enumerate}
	\item test(2,3)
	\item test(3+4,8)
	\item test(7,2*4)
	\item test(3+4,2*4)
\end{enumerate}

\subsubsection*{Svolgimento}
\begin{enumerate}
	\item CBV=CBN: entrambi terminano l'esecuzione in 2 passi; 
	\item CBV<CBN: CBV termina prima perch\`e valuta l'espressione 3+4 subito (una sola volta). CBN invece invoca prima la funzione, sostituendo ad x l'espressione. Essendo x presente 2 volte nella funzione, la valutazione viene effettuata 2 volte. 
	\item CBV>CBN: CBN termina prima perch\`e non valuta l'espressione 2*4 che non \`e utilizzata nella funzione test.
	\item CBV=CBN: entrambi terminano l'esecuzione in 4 passi; la differenza sta nel fatto che con CBV si svolgono prima le due operazioni (somma e moltiplicazione) e poi si chiama la funzione. Con CBN invece, si evita di calcolare 2*4, tuttavia si calcola 3+4 due volte.
\end{enumerate}



\subsection*{Esercizio 1.8}
Definire in Scala una funzione and che si comporta come il costrutto logico $x\&\&y$. 

\subsubsection*{Svolgimento}

\begin{lstlisting}
// and logico. b1 e b2 passati by-name 
def and(b1:=> Boolean, b2:=> Boolean): Boolean =
	if(!b1) false
	else if(b2) true
		 else false; 
\end{lstlisting}

Il parametro b1 viene passato by-name nella nostra soluzione. Tuttavia, esso pu\`o essere passato anche by-value visto che in ogni caso viene valutato. 











\section{I tipi semplici (note 3)}
\subsection*{Esempi di derivazione svolti}
 
\begin{enumerate}[label=\alph*)]
	\item $\emptyset\:\vdash$  if true then 5 + 7 else 2 :Nat
	\item $\emptyset\:\vdash$  $\fn$ x:T.x :T$\rightarrow$T
	\item $\emptyset\:\vdash$  ($\fn$ x:Bool.x) true :Bool
	\item f:Bool$\rightarrow$ Bool $\vdash$ f (if false then true else false) :Bool
	\item f:Bool$\rightarrow$ Bool $\vdash\:\fn$ x:Bool.f(if x then false else x) :Bool$\rightarrow$Bool
\end{enumerate}

\subsubsection*{Svolgimento}

\begin{enumerate}[label=\alph*), leftmargin=*]
	\item $\emptyset\:\vdash$  if true then 5 + 7 else 2 :Nat\\
		\scalebox{.75}{
		\parbox{1cm}{
    	\begin{prooftree}
			\AxiomC{$\checkmark$}
			\LeftLabel{\textsc{(TRUE)}}
			\UnaryInfC{$\emptyset\:\vdash$ true: Bool}
			\AxiomC{$\checkmark$}
			\LeftLabel{\textsc{(NAT)}}
			\UnaryInfC{$\emptyset\:\vdash$ 5 :Nat}
			\AxiomC{$\checkmark$}
			\LeftLabel{\textsc{(NAT)}}
			\UnaryInfC{$\emptyset\:\vdash$ 7 :Nat}
			\LeftLabel{\textsc{(SUM)}}
			\BinaryInfC{$\emptyset\:\vdash$ 5 + 7 :Nat}
			\AxiomC{$\checkmark$}
			\LeftLabel{\textsc{(NAT)}}
			\UnaryInfC{$\emptyset\:\vdash$  2 :Nat}
			\LeftLabel{\textsc{(IF-THEN-ELSE)}}
			\TrinaryInfC{$\emptyset\:\vdash$  if true then 5 + 7 else 2 :Nat}
		\end{prooftree}}}

    \vspace*{1 cm}

	\item $\emptyset\:\vdash$  $\fn$ x:T.x :T$\rightarrow$T  \\
		\scalebox{.75}{
		\parbox{1cm}{
		\begin{prooftree} 
			\AxiomC{x:T $\in$ x:T}
			\LeftLabel{\textsc{(VAR)}}
			\UnaryInfC{x:T $\vdash$  x :T}
			\LeftLabel{\textsc{(FUN)}}
			\UnaryInfC{$\emptyset\:\vdash$  $\fn$ x:T.x :T$\rightarrow$T}
		\end{prooftree}}
		}
	 
	 
    \vspace*{1 cm}
	\item $\emptyset\:\vdash$  ($\fn$ x:Bool.x) true :Bool\\
		\scalebox{.75}{
		\parbox{1cm}{
		\begin{prooftree} 
			\AxiomC{x:Bool $\in$ x:Bool}
			\LeftLabel{\textsc{(VAR)}}
			\UnaryInfC{x:Bool $\vdash$  x :Bool}
			\LeftLabel{\textsc{(FUN)}}
		    \UnaryInfC{$\emptyset\:\vdash\:\fn$ x:Bool.x :Bool$\rightarrow					$Bool}
			\AxiomC{$\checkmark$}
			\LeftLabel{\textsc{(TRUE)}}
			\UnaryInfC{$\emptyset\:\vdash$ true :Bool}
			\LeftLabel{\textsc{(APP)}}
			\BinaryInfC{$\emptyset\:\vdash$  ($\fn$ x:Bool.x) true :Bool}
		\end{prooftree}}
		}
		
	
    \vspace*{1 cm}	
	\item f:Bool$\rightarrow$ Bool $\vdash$ f (if false then true else false) :Bool\\
		Sia $\Gamma$ = f:Bool$\rightarrow$Bool \\
		\scalebox{.75}{
		\parbox{1cm}{
		\begin{prooftree}  
			\AxiomC{f:Bool$\rightarrow$ Bool $\in$ $\Gamma$}
			\LeftLabel{\textsc{(VAR)}}
			\UnaryInfC{$\Gamma$ $\vdash$ f :Bool $\rightarrow$ Bool} 
			\AxiomC{$\checkmark$}
			\LeftLabel{\textsc{(FALSE)}}
			\UnaryInfC{$\Gamma$ $\vdash$ false :Bool}
			\AxiomC{$\checkmark$}
			\LeftLabel{\textsc{(TRUE)}}
			\UnaryInfC{$\Gamma$ $\vdash$ true :Bool}
			\AxiomC{$\checkmark$}
			\LeftLabel{\textsc{(FALSE)}}
			\UnaryInfC{$\Gamma$ $\vdash$ false :Bool}
			\LeftLabel{\textsc{(IF-THEN-ELSE)}}
			\TrinaryInfC{$\Gamma$ $\vdash$ if false then true else false :Bool}
			\LeftLabel{\textsc{(APP)}}
			\BinaryInfC{$\Gamma$ $\vdash$ f (if false then true else 							false):Bool}
		\end{prooftree}}
		} 
	
	
    \vspace*{1 cm}
	\item f:Bool$\rightarrow$ Bool $\vdash\:\fn$ x:Bool.f(if x then false else x) :Bool$\rightarrow$Bool 
		Sia $\Gamma$ = f:Bool$\rightarrow$Bool \\
		\scalebox{.75}{
		\parbox{1cm}{
		\begin{prooftree}  
			\AxiomC{f:Bool $\rightarrow$ Bool $\in$ $\Gamma$, x:Bool}
			\LeftLabel{\textsc{(VAR)}}
			\UnaryInfC{$\Gamma$, x:Bool $\vdash$ f: Bool $\rightarrow$ Bool} 
			\LeftLabel{\textsc{(FUN)}}
			\UnaryInfC{$\Gamma$ $\vdash\:\fn$ x:Bool.f: ? $\rightarrow$ Bool 
			$\rightarrow$ Bool} 
			\AxiomC{x:Bool $\in$ $\Gamma$}
			\LeftLabel{\textsc{(VAR)}}
			\UnaryInfC{$\Gamma$ $\vdash$ x :?}
			\AxiomC{$\checkmark$}
			\LeftLabel{\textsc{(TRUE)}}
			\UnaryInfC{$\Gamma$ $\vdash$ true :Bool}
			\AxiomC{$\checkmark$}
			\LeftLabel{\textsc{(FALSE)}}
			\UnaryInfC{$\Gamma$ $\vdash$ false :Bool}
			\LeftLabel{\textsc{(IF-THEN-ELSE)}}
			\TrinaryInfC{$\Gamma$ $\vdash$ (if x then false else x):? 
			$\rightarrow$ Bool}
			\LeftLabel{\textsc{(APP)}}
			\BinaryInfC{$\Gamma\:\vdash\:\fn$ x:Bool.f(if x then false 
			else x)
			 :Bool$\rightarrow$Bool}
		\end{prooftree}}
		} \\
		Quindi ?=Bool.	
\end{enumerate}
 
    \vspace*{1 cm}

\subsection*{Esercizio 2.1}
Trovare un contesto $\Gamma$ tale che $\Gamma\:\vdash$ f x y : Bool sia derivabile.

\subsubsection*{Svolgimento}
 	\begin{prooftree} 
		\AxiomC{f:$T_2\:\rightarrow\:T_1\:\rightarrow$ Bool $\in\:\Gamma$}
		\LeftLabel{\textsc{(VAR)}}
		\UnaryInfC{$\Gamma\:\vdash$ f :$T_2\:\rightarrow\:T_1\:\rightarrow$
		 Bool}
		\AxiomC{x:$T_2\:\in$ x:$T_2$}
		\LeftLabel{\textsc{(VAR)}}
		\UnaryInfC{$\Gamma\:\vdash$ x :$T_2$}
		\LeftLabel{\textsc{(APP)}}
		\BinaryInfC{$\Gamma\:\vdash$ f x :$T_1\:\rightarrow$ Bool}
		\AxiomC{y :$T_1$ $\in$ $\Gamma$}
		\LeftLabel{\textsc{(VAR)}}
		\UnaryInfC{$\Gamma\:\vdash$ y :$T_1$}
		\LeftLabel{\textsc{(APP)}}
		\BinaryInfC{$\Gamma\:\vdash$ f x y :Bool}
	\end{prooftree} 
	
	Esiste un contesto per questo programma, in quanto il programma si pu\`o ottenere applicando le regole e costruendo l'albero di derivazione. \\
	Fanno parte del contesto le seguenti regole di tipo:
\begin{itemize}
	\item f:$T_2\:\rightarrow\:T_1\:\rightarrow$ Bool
	\item x :$T_2$
	\item y :$T_1$
\end{itemize}

\vspace*{2 cm}
Caso f (x y):\\
 	\begin{prooftree} 
		\AxiomC{f :$T_1\:\rightarrow$ Bool $\:\in\:\Gamma$}
		\LeftLabel{\textsc{(VAR)}}
		\UnaryInfC{$\Gamma\:\vdash$ f :$T_1\:\rightarrow$ Bool}
		\AxiomC{x:$T_2\:\rightarrow\:T_1\:\in\:\Gamma$}
		\LeftLabel{\textsc{(VAR)}}
		\UnaryInfC{$\Gamma\:\vdash$ x :$T_2\:\rightarrow\:T_1$}
		\AxiomC{y:$T_2\:\in\:\Gamma$}
		\LeftLabel{\textsc{(VAR)}}
		\UnaryInfC{$\Gamma\:\vdash$ y :$T_2$}
		\LeftLabel{\textsc{(APP)}}
		\BinaryInfC{$\Gamma\:\vdash$ x y :$T_1$}
		\LeftLabel{\textsc{(APP)}}
		\BinaryInfC{$\Gamma\:\vdash$ f (x y) :Bool}
	\end{prooftree} 
	
	Esiste un contesto anche per questo programma. \\
	Fanno parte del contesto le seguenti regole di tipo:
\begin{itemize}
	\item f:$T_1\:\rightarrow\:$ Bool
	\item x :$T_2\:\rightarrow\:T_1$
	\item y :$T_2$
\end{itemize}


\subsection*{Esercizio 2.2}
%Il giudizio $\Gamma\:\vdash$ x x : T \`e derivabile? Se s\`i, trovare una derivazione per qualche $\Gamma$, T, altrimenti provare che non \`e derivabile.

\subsubsection*{Svolgimento}
 	\begin{prooftree} 
		\AxiomC{$\Gamma\:\vdash$ x :$T_1\:\rightarrow\:$T}
		\AxiomC{$\Gamma\:\vdash$ x :$T_1$}
		\LeftLabel{\textsc{(APP)}}
		\BinaryInfC{$\Gamma\:\vdash$ x x :T}
	\end{prooftree} 

No, non \`e derivabile. Nel nostro sistema di tipi x non pu\`o assumere contemporaneamente il tipo T e il tipo $T_1\:\rightarrow\:$T. Risulterebbe derivabile se fosse presente il tipo ricorsivo nel sistema di tipi. \\






\newpage
\bibliographystyle{plain}
\bibliography{biblist}

\end{document} 
